% =============================================================================
%  Section 4.2 — Thực nghiệm benchmark
% =============================================================================

\section{Thực nghiệm benchmark}
\label{sec:benchmark_experiments}

\subsection{Triển khai hệ thống cho thực nghiệm}
\label{sec:benchmark_setup}

Thực nghiệm benchmark được thực hiện trên môi trường cục bộ với mục tiêu đo định lượng hiệu năng các thành phần của kiến trúc GraphRAG và so sánh với các hệ thống tham chiếu. Thiết kế thực nghiệm tuân thủ nguyên tắc kiểm soát chặt mọi yếu tố ngoài cơ chế truy xuất - từ phần cứng, mô hình ngôn ngữ, tham số tạo sinh đến tập câu hỏi - đều được giữ cố định để chênh lệch kết quả phản ánh đúng đặc trưng kiến trúc của từng hệ thống.

\textbf{Môi trường thực nghiệm.} Toàn bộ thực nghiệm được triển khai trên một máy cá nhân chạy Windows. Đồ thị tri thức Chèo được dựng trên Neo4j, pipeline GraphRAG với cả ba chiến lược tạo sinh được cài đặt đầy đủ và có thể chuyển đổi qua cấu hình. Mô hình ngôn ngữ lớn được sử dụng qua API để đảm bảo tính \textit{LLM-agnostic}, đồng thời một mô hình duy nhất được chọn cho toàn bộ thực nghiệm với tham số sinh cố định (temperature bằng 0, top-p bằng 1) - cách đặt này loại bỏ tính ngẫu nhiên trong sinh ngôn ngữ và đảm bảo các kết quả định lượng có khả năng tái lập.

\textbf{Hệ thống tham chiếu.} Một hệ thống tham chiếu được triển khai cùng môi trường để làm cơ sở cho thực nghiệm so sánh, đó là \textbf{RAG truyền thống} - đại diện cho cách tiếp cận RAG dựa trên tương đồng véc-tơ: tài liệu mô tả Chèo được chia thành các đoạn văn bản có độ dài cố định, nhúng véc-tơ bằng mô hình embedding và lưu trong vector store. Khi nhận câu hỏi, hệ thống lấy các đoạn văn có độ tương đồng cosine cao nhất làm ngữ cảnh trước khi gửi tới mô hình ngôn ngữ. Cả hai hệ thống GraphRAG và RAG truyền thống dùng cùng một mô hình ngôn ngữ và cùng tham số sinh để chênh lệch kết quả chỉ phản ánh khác biệt về cơ chế truy xuất.

\textbf{Quy mô mẫu thực nghiệm.} Mẫu thực nghiệm được phân chia theo mục đích từng thí nghiệm. Thí nghiệm so sánh GraphRAG với RAG truyền thống chạy trên toàn bộ 100 ca của CheoBench để đảm bảo phạm vi đánh giá đầy đủ. Các thí nghiệm phân tích thành phần - đo ảnh hưởng của mức độ truy xuất và định dạng ngữ cảnh đến hiệu năng tổng thể - chạy trên một tập con 20 ca được lấy mẫu phân tầng cân bằng theo ba nhóm Local, Community, Global. Tập con này được dùng chung cho mọi cấu hình thí nghiệm nhằm loại biến động giữa các tập câu hỏi và kiểm soát chi phí gọi mô hình ngôn ngữ.

\subsection{Độ đo đánh giá}

Hệ độ đo đánh giá của CheoBench gồm chín thành phần thuộc ba khối, phản ánh ba chiều cốt lõi của một hệ thống RAG: chất lượng tập thực thể được truy xuất, chất lượng ngữ cảnh đưa vào mô hình ngôn ngữ và chất lượng câu trả lời cuối cùng. Bốn độ đo IR truyền thống đo độ chính xác và độ bao phủ của danh sách thực thể truy xuất; ba độ đo ngữ nghĩa theo RAGAs~\cite{es2023ragas} đo chất lượng ngữ cảnh; hai độ đo còn lại theo RAGAs đo chất lượng câu trả lời. Chín thành phần được kết hợp thành ba điểm tổng hợp $S_{\text{retrieval}}$, $S_{\text{generation}}$ và $S_{\text{overall}}$, trong đó $S_{\text{overall}}$ đóng vai trò là chỉ số cuối cùng, quyết định chất lượng của cả hệ thống ở phần kết quả thực nghiệm.

\subsubsection{Chất lượng truy xuất}

Chất lượng truy xuất được đánh giá bởi bốn độ đo IR xác định kết hợp với ba độ đo ngữ nghĩa theo RAGAs. Nhóm IR đo độ chính xác và độ bao phủ của danh sách thực thể truy xuất so với tập thực thể mong muốn (ground truth):

\begin{itemize}
  \item \textbf{MAP}~\cite{manning2008ir} - trung bình độ chính xác (average precision) trên toàn bộ truy vấn; MAP thưởng các cấu hình đưa nhiều thực thể đúng lên vị trí đầu của danh sách truy xuất, phản ánh chất lượng xếp hạng tổng thể.

  \item \textbf{NDCG@10}~\cite{jarvelin2002ndcg} - đánh giá chất lượng xếp hạng ở mười vị trí đầu với phạt logarit theo vị trí; mức cắt $K = 10$ được chọn để phù hợp với kích thước ngữ cảnh điển hình mà pipeline đưa vào mô hình ngôn ngữ.

  \item \textbf{Precision} và \textbf{Recall} - tỷ lệ chính xác và độ bao phủ của kết quả truy xuất so với tập thực thể liên quan $R$ theo ground truth, tính trên toàn bộ danh sách truy xuất $L$ sau khi loại trùng lặp:
    \begin{equation}
      \text{Precision} = \frac{|L \cap R|}{|L|}, \qquad \text{Recall} = \frac{|L \cap R|}{|R|}.
    \end{equation}
\end{itemize}

Nhóm RAGAs bổ sung khía cạnh ngữ nghĩa, đo chất lượng ngữ cảnh đã được tổ chức trước khi đưa vào mô hình ngôn ngữ. Tất cả ba độ đo đều nằm trong thang $[0, 1]$. \textbf{Context Precision} và \textbf{Context Recall} được mô hình ngôn ngữ chấm điểm từ cặp (câu hỏi, ngữ cảnh) và (ground truth, ngữ cảnh) để ước lượng tỷ lệ thông tin liên quan và mức độ bao phủ tri thức cần thiết. \textbf{Context Entities Recall} được tính trực tiếp: với $E$ là tập các thực thể bắt buộc theo ground truth, độ đo là tỷ lệ thực thể trong $E$ xuất hiện trong ngữ cảnh được truy xuất:

\begin{equation}
\text{Context Entities Recall} = \frac{|E \cap \text{context}|}{|E|}.
\end{equation}

Điểm truy xuất tổng hợp $S_{\text{retrieval}}$ kết hợp bảy độ đo theo trọng số:

\begin{equation}
\begin{aligned}
S_{\text{retrieval}} =\;& 0,20 \cdot \text{MAP}
  + 0,15 \cdot \text{NDCG@10} \\
& + 0,10 \cdot \text{Precision}
  + 0,10 \cdot \text{Recall} \\
& + 0,20 \cdot \text{Context Precision}
  + 0,15 \cdot \text{Context Recall} \\
& + 0,10 \cdot \text{Context Entities Recall}.
\end{aligned}
\end{equation}

Khối IR nhận tổng trọng số $0,55$ - cao hơn khối RAGAs ($0,45$) - do các độ đo IR có thể tính toán trực tiếp và tái lập chính xác, trong khi RAGAs phụ thuộc vào phán đoán của mô hình ngôn ngữ và có biến động giữa các lần chạy. Trong từng khối, các độ đo bao quát chất lượng xếp hạng hoặc độ liên quan tổng thể (MAP, Context Precision) nhận trọng số cao nhất ($0,20$); các độ đo đo độ bao phủ ở phạm vi rộng (NDCG@10, Context Recall) nhận trọng số trung bình ($0,15$); các độ đo đơn khía cạnh hoặc phạm vi hẹp (Precision, Recall, Context Entities Recall) nhận trọng số thấp nhất ($0,10$).

\subsubsection{Chất lượng tạo sinh}

Chất lượng tạo sinh được đánh giá qua hai độ đo RAGAs theo cơ chế \textbf{LLM-as-judge} - mô hình ngôn ngữ đóng vai trò bên chấm điểm và trả về giá trị trong thang $[0, 1]$.

\textbf{Faithfulness} đo tỷ lệ nội dung câu trả lời thực sự được ngữ cảnh truy xuất hỗ trợ, được thiết kế để phát hiện hiện tượng ảo giác của mô hình. Giá trị độ đo được tính qua hai bước. Đầu tiên là \textit{tách mệnh đề}: mô hình ngôn ngữ liệt kê các khẳng định đơn lẻ trong câu trả lời thành tập $C = \{c_1, \ldots, c_n\}$. Sau đó là \textit{kiểm chứng từng mệnh đề}: với mỗi $c_i \in C$, mô hình được đưa ngữ cảnh truy xuất cùng mệnh đề và được hỏi mệnh đề đó có được hỗ trợ hay không; mệnh đề chỉ được tính là được hỗ trợ khi mô hình trả lời có. Faithfulness là tỷ lệ mệnh đề được hỗ trợ trên tổng số mệnh đề trích xuất:

\begin{equation}
\text{Faithfulness} = \frac{|\{c \in C : \text{được ngữ cảnh hỗ trợ}\}|}{|C|}.
\end{equation}

\textbf{Answer Relevance} được mô hình ngôn ngữ ước lượng trực tiếp từ cặp (câu hỏi, câu trả lời), cho điểm số liên tục phản ánh mức độ câu trả lời giải quyết đúng câu hỏi được đặt ra. Điểm tạo sinh $S_{\text{generation}}$ tổng hợp hai độ đo với trọng số bằng nhau:

\begin{equation}
S_{\text{generation}} = 0,5 \cdot \text{Faithfulness} + 0,5 \cdot \text{Answer Relevance}.
\end{equation}

\subsubsection{Điểm tổng hợp}

Điểm tổng hợp $S_{\text{overall}}$ kết hợp $S_{\text{retrieval}}$ và $S_{\text{generation}}$ với trọng số bằng nhau, đóng vai trò chỉ số chính trong phần kết quả thực nghiệm:

\begin{equation}
S_{\text{overall}} = 0,5 \cdot S_{\text{retrieval}} + 0,5 \cdot S_{\text{generation}}.
\end{equation}

\subsection{Kết quả thực nghiệm}
\label{sec:benchmark_results}

Phần kết quả lần lượt báo cáo ba thí nghiệm tương ứng ba mục tiêu. \textbf{Mục tiêu thứ nhất (MT1)} so sánh tổng thể GraphRAG với RAG truyền thống trên toàn bộ 100 ca của CheoBench để xác định lợi ích của kiến trúc đề xuất. Hai mục tiêu còn lại phân tích từng thành phần của GraphRAG trên tập con 20 ca: \textbf{mục tiêu thứ hai (MT2)} đo ảnh hưởng của mức độ truy xuất trong đồ thị tri thức đến hiệu năng tổng thể, \textbf{mục tiêu thứ ba (MT3)} đo tác động của định dạng biểu diễn ngữ cảnh đến khả năng suy luận của mô hình ngôn ngữ.

\subsubsection{So sánh với hệ thống tham chiếu (MT1)}
\label{sec:results_baseline}

Kết quả so sánh GraphRAG với RAG truyền thống trên 100 ca của CheoBench (Bảng~\ref{tab:baseline_results}) cho ba phát hiện đáng chú ý, lần lượt theo ba khối độ đo IR, Context và Generation.

Ở khối IR thuần, hai hệ thống có hiệu năng tương đương trên các độ đo xếp hạng - MAP đạt $0,500$ cho GraphRAG so với $0,520$ cho RAG, NDCG@10 đạt $0,579$ so với $0,588$ - thậm chí RAG nhỉnh hơn nhẹ ở Precision tuyệt đối ($0,273$ so với $0,173$). Khác biệt này không phản ánh chất lượng xếp hạng mà phản ánh cơ chế trả về các thực thể ứng viên: RAG chỉ lấy các đoạn văn bản có độ tương đồng cosine cao nhất nên danh sách ngắn và đặc, GraphRAG mở rộng truy xuất qua cấu trúc đồ thị nên trả về nhiều thực thể ứng viên hơn - hệ quả là Recall của GraphRAG vượt RAG ($0,707$ so với $0,598$) còn Precision lại thấp hơn khá nhiều. Sự tương đương ở khối IR do đó cho thấy việc tìm đúng thực thể hay đoạn văn bản liên quan đến câu hỏi không phải là khác biệt cốt lõi giữa hai kiến trúc.

Khoảng cách thực sự bộc lộ ở khối Context. Context Recall đạt $0,964$ cho GraphRAG so với $0,418$ cho RAG - chênh lệch tới $+0,546$ điểm. Nguyên nhân nằm ở cách hai hệ thống tổ chức ngữ cảnh trước khi đưa vào mô hình ngôn ngữ: RAG cắt tài liệu mô tả Chèo thành các đoạn cố định rồi lấy theo độ tương đồng véc-tơ, dẫn đến tình trạng thông tin về một thực thể có thể nằm rải rác ở nhiều đoạn nhưng chỉ một số đoạn được lấy ra; trong khi GraphRAG truy xuất theo cấu trúc đồ thị, đưa cả thực thể trung tâm lẫn các thực thể có quan hệ trực tiếp vào ngữ cảnh một cách hệ thống (Mục~\ref{sec:multi_granularity}). Context Precision của GraphRAG ($0,867$) và Context Entities Recall ($0,982$) cũng vượt xa RAG ($0,513$ và $0,637$), khẳng định ngữ cảnh từ đồ thị tri thức vừa liên quan hơn vừa đầy đủ hơn so với ngữ cảnh từ vector store.

Chất lượng ngữ cảnh khác biệt dẫn đến chênh lệch rõ rệt ở khối Generation. Faithfulness của GraphRAG đạt $0,986$ - gần như mọi mệnh đề trong câu trả lời đều có cơ sở trong ngữ cảnh - so với $0,734$ của RAG, tức RAG có khoảng $27\%$ mệnh đề không được ngữ cảnh hỗ trợ và phải dựa vào tri thức nội tại của mô hình ngôn ngữ. Kết quả này phù hợp với hai cơ chế chống ảo giác đã thiết kế trong G-Retrieval: Danh mục thực thể (Mục~\ref{sec:entity_catalog}) ràng buộc tên gọi không vượt phạm vi đồ thị, và ngữ cảnh có cấu trúc rõ ràng giúp mô hình ngôn ngữ phân biệt được thực thể nào là trung tâm câu hỏi. Answer Relevance ($0,974$ so với $0,630$) là hệ quả của hai yếu tố trên: khi ngữ cảnh thiếu thông tin then chốt (Context Recall $0,418$ của RAG), câu trả lời lan man hoặc lệch hướng so với câu hỏi đặt ra.

Tổng hợp lại, lợi ích của GraphRAG so với RAG truyền thống chuyển thẳng vào điểm tổng hợp: $S_{\text{retrieval}}$ chênh $+0,183$ điểm chủ yếu từ khối Context, $S_{\text{generation}}$ chênh $+0,298$ điểm trực tiếp từ Faithfulness và Answer Relevance, và $S_{\text{overall}}$ cuối cùng đạt $0,835$ cho GraphRAG so với $0,595$ cho RAG - chênh $+0,240$ điểm trên thang $[0, 1]$. Khoảng cách này lớn ở mức kiến trúc chứng minh rằng: cấu trúc đồ thị tri thức không chỉ cải thiện chất lượng câu trả lời ở mức tinh chỉnh mà còn thay đổi căn bản chất lượng ngữ cảnh được cung cấp cho các mô hình ngôn ngữ.

\begin{table}[h]
\centering
\caption{So sánh GraphRAG với RAG truyền thống trên CheoBench (n=100 ca).}
\label{tab:baseline_results}
\begin{tabular}{llcc}
\hline
\textbf{Khối} & \textbf{Độ đo} & \textbf{GraphRAG} & \textbf{RAG} \\
\hline
\multirow{4}{*}{IR}
& MAP            & 0,500 & 0,520 \\
& NDCG@10        & 0,579 & 0,588 \\
& Precision      & 0,173 & 0,273 \\
& Recall         & 0,707 & 0,598 \\
\hline
\multirow{3}{*}{Context (RAGAs)}
& Context Precision         & \textbf{0,867} & 0,513 \\
& Context Recall            & \textbf{0,964} & 0,418 \\
& Context Entities Recall   & \textbf{0,982} & 0,637 \\
\hline
\multirow{2}{*}{Generation (RAGAs)}
& Faithfulness     & \textbf{0,986} & 0,734 \\
& Answer Relevance & \textbf{0,974} & 0,630 \\
\hline
\multirow{3}{*}{Tổng hợp}
& $S_{\text{retrieval}}$  & \textbf{0,691} & 0,508 \\
& $S_{\text{generation}}$ & \textbf{0,980} & 0,682 \\
& $S_{\text{overall}}$    & \textbf{0,835} & 0,595 \\
\hline
\end{tabular}
\end{table}

\subsubsection{Ảnh hưởng của mức độ truy xuất (MT2)}
\label{sec:results_granularity}

Hiệu năng của GraphRAG được đo trên cùng tập 20 ca khi cấu hình chỉ dùng một mức truy xuất duy nhất, sau đó so sánh với cấu hình kết hợp cả bốn mức (Bảng~\ref{tab:granularity_results}). Trong các cấu hình đơn lẻ, Triplet đạt điểm cao nhất ($0,672$) do đáp ứng tốt phần lớn truy vấn vốn chỉ cần quan hệ trực tiếp giữa hai thực thể, với chi phí ngữ cảnh thấp. Path và Subgraph có điểm thấp hơn do trả về thêm nhiều quan hệ trung gian không thiết yếu, làm pha loãng tín hiệu. Cấu hình kết hợp đa mức độ đạt $0,811$ điểm, vượt mức tốt nhất trong các cấu hình đơn lẻ tới $+0,139$ điểm - xác nhận bốn mức bổ sung cho nhau thay vì cạnh tranh nhau.

\begin{table}[h]
\centering
\caption{Ảnh hưởng của mức độ truy xuất tới hiệu năng hệ thống (n=20 ca).}
\label{tab:granularity_results}
\begin{tabular}{lccc}
\hline
\textbf{Mức độ truy xuất} & \textbf{$S_{\text{retrieval}}$} & \textbf{$S_{\text{generation}}$} & \textbf{$S_{\text{overall}}$} \\
\hline
Node              & 0,382 & 0,917 & 0,649 \\
Triplet           & 0,383 & 0,962 & 0,672 \\
Path              & 0,322 & 0,835 & 0,578 \\
Subgraph          & 0,284 & 0,971 & 0,628 \\
\textbf{Kết hợp đa mức độ} & \textbf{0,634} & \textbf{0,989} & \textbf{0,811} \\
\hline
\end{tabular}
\end{table}

\subsubsection{Ảnh hưởng của định dạng biểu diễn ngữ cảnh (MT3)}
\label{sec:results_format}

Hiệu năng của GraphRAG được đo trên cùng tập 20 ca khi cấu hình chỉ dùng một định dạng biểu diễn ngữ cảnh duy nhất, sau đó so sánh với cấu hình kết hợp cả bốn định dạng (Bảng~\ref{tab:format_results}). Trong các định dạng đơn lẻ, JSON có cấu trúc đạt điểm cao nhất với $0,806$, tiếp theo là văn bản tự nhiên $0,761$ và bảng kề $0,750$; chuỗi nút tuyến tính thấp nhất ở $0,649$. Khoảng cách giữa định dạng tốt nhất và kém nhất lên tới $+0,157$ điểm, khẳng định kết quả của mô hình ngôn ngữ bị ảnh hưởng lớn bởi cách trình bày tri thức chứ không chỉ bởi mỗi lượng tri thức được cung cấp. Cấu hình kết hợp nhiều định dạng đạt $0,803$ - tương đương JSON đơn lẻ chứ không vượt lên, phản ánh hiện tượng lợi ích biên giảm dần khi tăng kích thước ngữ cảnh: khi định dạng tốt nhất đã cung cấp đủ tín hiệu, việc bổ sung thêm các bản trình bày khác chỉ tăng độ dài ngữ cảnh mà không mang lại tri thức mới.

\begin{table}[h]
\centering
\caption{Ảnh hưởng của định dạng biểu diễn ngữ cảnh tới hiệu năng hệ thống (n=20 ca).}
\label{tab:format_results}
\begin{tabular}{lccc}
\hline
\textbf{Định dạng} & \textbf{$S_{\text{retrieval}}$} & \textbf{$S_{\text{generation}}$} & \textbf{$S_{\text{overall}}$} \\
\hline
Văn bản tự nhiên        & 0,567 & 0,955 & 0,761 \\
\textbf{Cấu trúc JSON}  & \textbf{0,626} & \textbf{0,985} & \textbf{0,806} \\
Bảng kề                 & 0,530 & 0,970 & 0,750 \\
Chuỗi nút               & 0,560 & 0,738 & 0,649 \\
Kết hợp nhiều định dạng & 0,631 & 0,976 & 0,803 \\
\hline
\end{tabular}
\end{table}

\subsection{Phân tích và thảo luận}
\label{sec:benchmark_discussion}

Các kết quả thực nghiệm benchmark cho phép rút ra một số nhận định tổng hợp về kiến trúc GraphRAG, đồng thời gợi mở các hàm ý phương pháp luận đối với việc đánh giá các hệ thống RAG dựa trên đồ thị tri thức.

Sự khác biệt giữa GraphRAG và RAG truyền thống ở miền tri thức Chèo cho thấy vai trò của cấu trúc đồ thị không nằm ở pha truy xuất mà nằm ở pha tổ chức ngữ cảnh. Trên 100 ca của CheoBench, hai hệ thống cho hiệu năng IR thuần tương đương ở các độ đo xếp hạng (chênh lệch không quá $0,02$ điểm ở MAP và NDCG@10), nhưng tách biệt rõ rệt ở khối Context và khối Generation - đặc biệt là Context Recall ($0,964$ so với $0,418$) và Faithfulness ($0,986$ so với $0,734$). Lợi ích của việc khai thác đồ thị tri thức do đó không nằm ở khả năng \textit{tìm đúng} thực thể có liên quan - vốn dĩ độ tương đồng véc-tơ thực hiện đủ tốt trong một miền có cấu trúc tương đối rõ - mà nằm ở khả năng \textit{tổ chức} các thực thể đó thành ngữ cảnh có liên kết, đầy đủ và có ràng buộc về sự tồn tại trước khi đưa vào mô hình ngôn ngữ. Điều đó chứng minh một điều rằng, một phương pháp truy xuất thông tin tốt, cần đi cùng một phương pháp tổ chức ngữ cảnh tốt mới có thể đem lại hiệu quả tối đa khi các mô hình ngôn ngữ sinh câu trả lời.

Bên trong phân hệ G-Retrieval, các thành phần thể hiện hai kiểu tương tác trái ngược nhau. Bốn mức độ truy xuất Node, Triplet, Path và Subgraph tạo \textit{hiệu ứng cộng hưởng} khi được kết hợp - cấu hình đa mức vượt cấu hình đơn lẻ với nhiều nhất $+0,139$ điểm, lớn hơn đáng kể so với khoảng cách giữa các mức đơn lẻ với nhau - cho thấy bốn mức đo các khía cạnh không thể thay thế lẫn nhau của cấu trúc đồ thị, và cơ chế định tuyến đa mức theo loại truy vấn (Mục~\ref{sec:multi_granularity}) là một thành phần thiết kế thiết yếu chứ không phải tuỳ chọn tăng cường. Ngược lại, bốn định dạng biểu diễn ngữ cảnh thể hiện hiện tượng \textit{lợi ích biên giảm dần}: cấu hình kết hợp ($0,803$) gần như ngang với JSON đơn lẻ ($0,806$) trong khi khoảng cách giữa định dạng tốt nhất và kém nhất lên tới $+0,157$ điểm. Sự đối lập giữa hai pattern này thể hiện rằng: đa dạng hoá nguồn tri thức đem lại lợi ích cộng dồn, còn đa dạng hoá cách trình bày cùng một nguồn tri thức không cộng dồn được, một ngữ cảnh đủ tốt ổn định hơn so với đa dạng ngữ cảnh làm rối loạn thông tin khi cung cấp chúng cho các mô hình ngôn ngữ lớn.

Bên cạnh các phát hiện về kiến trúc, kết quả thí nghiệm so sánh với hệ thống tham chiếu (Mục~\ref{sec:results_baseline}) đồng thời minh chứng cho tính phù hợp của bộ benchmark CheoBench như một công cụ đo lường. Trên $100$ cặp chấm chung, GraphRAG vượt RAG ở $94/100$ câu mà không có câu nào hai hệ thống đồng điểm; khoảng cách điểm tổng hợp trung bình $0,24$ đi kèm phân phối trải rộng - $72$ câu chênh ít nhất $0,10$ và $33$ câu chênh ít nhất $0,30$. Phổ chênh lệch ổn định và rộng cho thấy CheoBench phân biệt được hai hệ thống ở nhiều cấp độ chất lượng chứ không dồn về hai cực, đáp ứng yêu cầu cơ bản của một bộ benchmark phân biệt theo nghĩa cổ điển.